\documentclass[11pt]{article}
\usepackage[margin=1in]{geometry}
\usepackage{amsmath, amssymb, amsthm, bm, mathtools}
\usepackage{booktabs}
\usepackage{enumitem}
\usepackage{microtype}
\usepackage{hyperref}
\hypersetup{colorlinks=true, linkcolor=blue, urlcolor=blue, citecolor=blue}

\title{Matrices and Their Operations}
\author{ECON 320 -- Introduction to Mathematical Economics }
\date{Fall 2025}

\newcommand{\R}{\mathbb{R}}

\begin{document}
\maketitle



\section{Introduction}
A \textbf{matrix} is a rectangular array of numbers arranged in rows and columns. 
They provide a compact representation of systems of linear equations, econometric models, and linear transformations.
Matrices are widely used in economics: equilibrium analysis, input--output models, comparative statics, and econometric estimation.



\section{Types of Matrices}
\begin{itemize}
    \item \textbf{Row Matrix:} One row only. Example: $[\,3 \; 5 \; 7\,]$.
    \item \textbf{Column Matrix:} One column only. Example: $\begin{bmatrix} 2 \\ 4 \\ 6 \end{bmatrix}$.
    \item \textbf{Square Matrix:} Same number of rows and columns ($n\times n$).
    \item \textbf{Diagonal Matrix:} All off-diagonal entries are zero.
    \item \textbf{Identity Matrix $I$:} Diagonal matrix with $1$’s on the diagonal. Acts like $1$ in multiplication.
    \item \textbf{Zero Matrix $0$:} All entries are zero.
    \item \textbf{Symmetric Matrix:} $A = A'$.
\end{itemize}



\section{Basic Operations}
\subsection{Addition and Subtraction}
Defined element-wise, only possible if matrices have the same dimensions.

\subsection{Scalar Multiplication}
Multiply every entry of a matrix by the scalar.

\subsection{Matrix Multiplication}
If $A$ is $m \times n$ and $B$ is $n \times p$, then $AB$ is $m \times p$:
\[
(AB)_{ij} = \sum_{k=1}^n a_{ik} b_{kj}.
\]

\subsection{Transpose}
The transpose $A'$ (or $A^T$) is obtained by flipping rows into columns.  
Properties:
\[
(A')' = A, \quad (A+B)' = A' + B', \quad (AB)' = B'A'.
\]



\section{Determinants}
The determinant is a scalar value associated with square matrices.  
\subsection{Definition}
For a $2 \times 2$ matrix:
\[
\det\begin{bmatrix} a & b \\ c & d \end{bmatrix} = ad - bc.
\]
For higher-order matrices, expansion by minors or Laplace expansion is used.

\subsection{Properties}
\begin{itemize}
    \item $\det(AB) = \det(A)\det(B)$.
    \item $\det(A') = \det(A)$.
    \item If $\det(A) = 0$, $A$ is \textbf{singular} and not invertible.
\end{itemize}



\section{Inverse of a Matrix}
The inverse $A^{-1}$ (if it exists) satisfies:
\[
AA^{-1} = A^{-1}A = I.
\]
For a $2 \times 2$ matrix:
\[
A^{-1} = \frac{1}{ad - bc}\begin{bmatrix} d & -b \\ -c & a \end{bmatrix}.
\]



\section{Jacobian Matrix}
The \textbf{Jacobian} of a vector-valued function $F:\R^n \to \R^m$ is the $m \times n$ matrix of first-order partial derivatives:
\[
J(x) = 
\begin{bmatrix}
\frac{\partial f_1}{\partial x_1} & \cdots & \frac{\partial f_1}{\partial x_n} \\
\vdots & \ddots & \vdots \\
\frac{\partial f_m}{\partial x_1} & \cdots & \frac{\partial f_m}{\partial x_n}
\end{bmatrix}.
\]
In economics, Jacobians are used in comparative statics, general equilibrium, and optimization problems.




\section*{Examples}

\subsection{Determinant Example}
Let 
\[
A = \begin{bmatrix} 2 & 1 \\ 3 & 4 \end{bmatrix}.
\]
Compute $\det(A)$:
\[
\det(A) = (2)(4) - (1)(3) = 8 - 3 = 5.
\]
Since $\det(A) \neq 0$, $A$ is invertible.



\subsection{Jacobian Example}
Consider the demand system:
\[
q_1(p_1, p_2) = 10 - 2p_1 + p_2, 
\quad q_2(p_1, p_2) = 20 + p_1 - 3p_2.
\]
The Jacobian is:
\[
J = 
\begin{bmatrix}
\frac{\partial q_1}{\partial p_1} & \frac{\partial q_1}{\partial p_2} \\
\frac{\partial q_2}{\partial p_1} & \frac{\partial q_2}{\partial p_2}
\end{bmatrix}
=
\begin{bmatrix}
-2 & 1 \\
1 & -3
\end{bmatrix}.
\]







\end{document}
